Las viñas orientadas a la exportación enfrentan un problema de planificación del envasado en el que cada producto terminado corresponde a una combinación única de vino y etiqueta, determinada por el idioma, la normativa sanitaria y los requerimientos comerciales de cada mercado de destino. Dado que los pronósticos de los importadores son poco precisos, comprometer anticipadamente el etiquetado expone a la viña a mantener inventario en combinaciones que la demanda no valida, mientras que producir exclusivamente contra pedido multiplica los tiempos de preparación y presiona la capacidad de la línea. La postergación del etiquetado, consistente en almacenar botellas llenas sin etiquetar hasta disponer de información de demanda más precisa, constituye un punto intermedio cuyo valor económico es necesario cuantificar.

Este informe corresponde a la primera etapa del proyecto, orientada a caracterizar el problema y a definir la estrategia de modelación. Se realizó un análisis exploratorio de la instancia base, correspondiente a una viña con dos vinos, cinco combinaciones vino--etiqueta y una demanda representada mediante un árbol de escenarios de tres períodos y once nodos. El análisis verificó la consistencia probabilística del árbol, caracterizó la demanda esperada y su variabilidad por producto, y estimó la estructura de correlación entre productos. Los resultados muestran correlaciones heterogéneas entre etiquetas que comparten un mismo vino base, condición bajo la cual la literatura reporta los mayores beneficios de la postergación, y sugieren que la capacidad de la línea no constituiría la restricción activa del problema.

Sobre esa base se selecciona la optimización estocástica multietapa como metodología y se formula un modelo preliminar de embotellado y etiquetado, junto con los indicadores de desempeño y el diseño experimental que permitirán cuantificar el valor de la postergación en las etapas siguientes.

\vfill
\noindent {\bf Palabras Claves}: postergación del etiquetado, optimización estocástica multietapa, planificación del envasado, incertidumbre en la demanda, industria vitivinícola.